\documentclass[10pt,a4paper]{amsart}
\usepackage[margin=19mm]{geometry}
\usepackage{amsmath,amssymb,mathtools}
\usepackage[scr=rsfs,cal=euler]{mathalfa}
\usepackage{xfrac}
\usepackage[unicode,hidelinks,bookmarks=false]{hyperref}
\newcommand{\ie}{i.e.\ }
\newcommand{\cf}{cf.\ }
\newcommand{\resp}{respectively\ }
\newcommand{\Subsection}{\subsection}
\providecommand{\nobreakdash}{\nobreak\hbox{-}}
\setlength{\emergencystretch}{2em}
\multlinegap0pt
\usepackage{amssymb}
\usepackage{algorithm}\usepackage{algpseudocode}
\usepackage{xcolor}
\newtheorem{definition}{Definition}
\newtheorem{proposition}{Proposition}
\newcommand{\Ineg}{I^{-}}
\newcommand{\Ipos}{I^{+}}
\newcommand{\Izero}{I^{0}}
\newcommand{\card}{\operatorname{card}}
\title{Accelerating Fourier--Motzkin elimination:\\ redundancy removal and the choice of variable elimination order}
\author{Shashaank Khanna}
\date{Source: arXiv:2609.07960v1 (CC BY 4.0)}
\begin{document}
\maketitle

\noindent\emph{Source and licence.} Accelerating Fourier--Motzkin elimination:\\ redundancy removal and the choice of variable elimination order. Version arXiv:2609.07960v1 (CC BY 4.0), distributed by arXiv under the Creative Commons Attribution 4.0 International licence, http://creativecommons.org/licenses/by/4.0/, as recorded in the arXiv licence metadata for this exact version and shown on the abstract page https://arxiv.org/abs/2609.07960v1. Full licence notice: \texttt{licenses/arxiv-2609.07960-CC-BY-4.0.txt}.\par\smallskip

\noindent\textit{Editorial scope.} Counted unit: the article's proposition prop:incompat (source0.tex lines 220--223). The article's complete original proof of it is delivered with it (source0.tex lines 225--251, one \texttt{proof} environment of 27 lines). No mathematics is written, repaired or re-derived: the statement and the proof are the article's own lines, in the article's own notation, and no retained line is substituted or re-flowed.  Retained uncounted as the source-local context the proof needs: lines 196--199; lines 206--209.  Nothing was omitted from any retained span, so the package carries no omission record.  No retained line contains a citation, so the wrapper carries no bibliography and no bibliography entry.  No retained line contains a figure environment, an image-inclusion command or drawing code, so no image is delivered and the package declares no asset dependency.  Editorial formatting only, in addition: an AMS \texttt{amsart} preamble that restates the article's own declarations and macros used by the retained lines, namely the environments \texttt{definition}, \texttt{proposition} on separate counters, together with the standard packages those lines need.  The retained environments print in this document as Definition 1, Proposition 1 in order of appearance, whereas the article prints the same items with the numbering of its own full document.  The complete native source of the article as distributed in the arXiv e-print for this exact version is bundled unchanged as \texttt{sources/209679/source0.tex}.
\begin{definition}
\label{def:imbert}
The \emph{history} $H_i$ of an inequality $i$ is the set of root inequalities used in its derivation: the history of a root inequality is the inequality itself, a copied inequality inherits its history, and the history of an inequality formed from a pair is the union of the histories of the two members of the pair. The \emph{explicit variables} $E_i$ are the variables that were eliminated by forming pairs somewhere in the derivation of $i$. The \emph{implicit variables} $I_i$ are the variables that occur in at least one member of $H_i$, but do not occur in $i$, and are not in $E_i$. Finally $O_k$ denotes the set of variables eliminated in the first $k$ steps.
\end{definition}

\begin{equation}
\card(H_i)\ \le\ 1+\card\!\big(E_i\cup(I_i\cap O_k)\big),
\label{eq:imbertcond}
\end{equation}

\begin{proposition}
\label{prop:incompat}
There is a system of linear inequalities and an elimination order for which the following procedure outputs a system whose solution set strictly contains the projection: after each elimination step, delete the inequalities that violate \eqref{eq:imbertcond}, then delete the inequalities that the LP test finds redundant, and carry the records of Definition~\ref{def:imbert} over to the next step. In the example below the procedure outputs the empty system.
\end{proposition}

\begin{proof}
Consider the following system in the variables $(w,x,y,z)$, which are otherwise unconstrained:
\begin{equation}
\begin{aligned}
(\mathrm{E}1)\quad -x+w &\le 0, &\qquad (\mathrm{E}2)\quad x+y &\le 0,\\
(\mathrm{E}3)\quad -x+y+z-w &\le 0, &\qquad (\mathrm{E}4)\quad x-y-z+w &\le 0.
\end{aligned}
\label{eq:counterexample}
\end{equation}
We eliminate $x$ and then $z$. The projection onto $(w,y)$ is the half-plane $\{y+w\le 0\}$. Indeed, adding (E1) and (E2) shows that $y+w\le 0$ holds on the solution set, and conversely, if $y+w\le 0$ then $x=w$ and $z=2w-y$ satisfy all four inequalities.

\emph{Elimination of $x$.} Here $\Ipos=\{\mathrm{E}2,\mathrm{E}4\}$, $\Ineg=\{\mathrm{E}1,\mathrm{E}3\}$ and $\Izero=\emptyset$, and the four pairs give
\begin{equation}
\begin{aligned}
(\mathrm{F}1)\quad y+w&\le 0, & H&=\{\mathrm{E}1,\mathrm{E}2\},\\
(\mathrm{F}2)\quad 2y+z-w&\le 0, & H&=\{\mathrm{E}2,\mathrm{E}3\},\\
(\mathrm{F}3)\quad -y-z+2w&\le 0, & H&=\{\mathrm{E}1,\mathrm{E}4\},\\
(\mathrm{F}4)\quad 0&\le 0, & H&=\{\mathrm{E}3,\mathrm{E}4\},
\end{aligned}
\label{eq:afterx}
\end{equation}
each with $E=\{x\}$. Imbert's test flags nothing, since $\card(H)=2\le 1+\card(\{x\})$ for each of the four. The LP test, on the other hand, finds (F1) redundant, because the sum of (F2) and (F3) is $y+w\le 0$, and it finds (F4) redundant trivially.

\emph{Imbert's test alone, LP test at the end.} Suppose we keep all four inequalities of \eqref{eq:afterx} and eliminate $z$. Then (F2) and (F3) form a pair, giving a second copy of $y+w\le 0$ with $H=\{\mathrm{E}1,\mathrm{E}2,\mathrm{E}3,\mathrm{E}4\}$ and $E=\{x,z\}$, while (F1) and (F4) are copied. Condition \eqref{eq:imbertcond} fails for the new copy, since $4\not\le 1+2$, and Imbert's test deletes it. This is correct: the new copy is redundant given the copy (F1). A final LP pass then leaves $\{y+w\le 0\}$, which is the projection.

\emph{The interleaved procedure.} Now suppose that after the elimination of $x$ we delete (F1) and (F4), as the LP test instructs, leaving (F2) and (F3). Eliminating $z$ produces the single inequality $y+w\le 0$ with $H=\{\mathrm{E}1,\mathrm{E}2,\mathrm{E}3,\mathrm{E}4\}$ and $E=\{x,z\}$. Condition \eqref{eq:imbertcond} fails, the inequality is deleted, and the output is the empty system, whose solution set is $\mathbb{R}^2$.
\end{proof}

\end{document}
